% !TEX root = report.tex

\section{Summary}
\label{sec:conclusion-summary}

Carbon-aware scheduling depends on carbon-intensity forecasts to choose when and where to run flexible workloads.
The existing forecasting literature treats this as an accuracy problem and optimizes its models against MAPE, but the scheduler's decision is fundamentally an ordering problem: given a window of candidate hours or regions, the scheduler picks whichever has the lowest predicted emissions.
This thesis argues that the dimension of forecast quality that actually drives scheduling outcomes is rank-preservation, formalized as the concordance index $C(\hat{f}, f^{\text{act}})$, and not pointwise accuracy.
To demonstrate this we built LiteCast, a lightweight per-source SARIMAX forecaster trained on a month-long rolling window of hourly energy, weather, and demand data, and evaluated it against the state-of-the-art deep-learning forecaster CarbonCast across $48$ regions and a real cluster trace.

\section{Contributions}
\label{sec:conclusion-contributions}

This thesis makes three contributions.

\textit{Analytic.}
We show that the scheduler's chosen start hour (continuous jobs) or chosen hour set (interruptible jobs) depends on the forecast only through the ordering it induces over the window.
Two forecasts with identical MAPE can yield different schedules whenever their induced orderings differ, and two forecasts with very different MAPE can yield identical schedules whenever their orderings agree on the relevant comparisons.
The concordance index is the natural quantification of this rank-agreement, and is therefore the forecast-quality metric most directly aligned with scheduling outcomes.

\textit{Systems.}
We design and implement LiteCast, a per-source SARIMAX pipeline that produces hourly carbon-intensity forecasts spanning more than a week from a month of recent energy, weather, and demand observations.
LiteCast trains in seconds on a single CPU, refits on each job submission, and requires no GPU infrastructure or multi-year historical data.

\textit{Empirical.}
Across $48$ worldwide regions, LiteCast outperforms CarbonCast on both concordance and MAPE, and the gap on realized carbon savings tracks concordance rather than MAPE.
The result holds across continuous and interruptible job models, across slack lengths from $24$ to $168$ hours, across a spatial-migration evaluation over all grid pairs, and on the KIT cluster-trace replay where job lengths and slacks are drawn from real user behavior rather than from a synthetic sweep.

\section{Future Work}
\label{sec:conclusion-future-work}

Two directions follow naturally from this work.

The first is to make concordance a training objective rather than only an evaluation metric.
LiteCast achieves high concordance as a consequence of its short training window and per-source decomposition, but its underlying SARIMAX fit is still performed against a pointwise likelihood.
A forecaster that optimizes a rank-aware loss directly would be a more principled realization of the thesis's central argument and may close more of the remaining gap to the oracle.

The second is to extend LiteCast beyond a single base model.
SARIMAX captures the linear, exogenously-driven structure of most grids well but is not uniformly best across the $48$ regions in our set: regions whose carbon intensity is dominated by non-linear or weather-coupled dynamics may benefit from alternative base models such as gradient-boosted trees or compact recurrent networks.
A natural next step is to evaluate a family of lightweight forecasters under the same concordance-first design and let the deployment select among them per region, providing alternatives where SARIMAX performs worst while preserving the lightweight footprint that motivates the rest of the design.
